\chapter{Working with GitHub}
\label{ch:github}

\chapabstract{\GitHub{} hosts \Git{} repositories and adds collaboration
on top: pull requests, issues, and access control. This chapter covers
the workflow you need to contribute to a shared repository.}

\section{Remotes}

A remote is a named reference to another copy of the repository. Cloning
a repository creates one automatically, named \code{origin}.

\begin{shcode}
git remote -v                                # list remotes and their URLs
git remote add origin git@github.com:user/repo.git
git push -u origin main   # push, and remember this as the default upstream
git push                  # afterwards, this alone is enough
git fetch origin          # download commits, do not touch the working dir
git pull                  # fetch, then merge into the current branch
\end{shcode}

\begin{tipbox}{SSH over HTTPS}
An HTTPS remote (\code{https://github.com/...}) asks for a username and
a personal access token on every push. An SSH remote
(\code{git@github.com:...}) uses a key pair once configured
(\code{ssh-keygen}, then add the public key under \GitHub{} \(\to\)
Settings \(\to\) SSH and GPG keys) and never prompts again. Prefer SSH
for any repository you push to regularly.
\end{tipbox}

\section{Forking versus branching}

Two different ways to propose a change, depending on your access level.

\begin{tabularx}{\linewidth}{@{}l>{\raggedright\arraybackslash}X@{}}
\toprule
\textbf{You have write access} & \textbf{You do not}\\
\midrule
Create a branch directly in the shared repository. &
  \textbf{Fork} the repository into your own account (a full server-side
  copy), clone your fork, and work there.\\
\bottomrule
\end{tabularx}

\begin{shcode}
# after forking user/repo into your-name/repo on github.com
git clone git@github.com:your-name/repo.git
cd repo
git remote add upstream git@github.com:user/repo.git   # the original
git fetch upstream
git merge upstream/main        # keep your fork up to date
\end{shcode}

\section{Pull requests}

A \textbf{pull request} (PR) asks the maintainers of a repository to
merge a branch (yours, or one in your fork) into theirs. It is where
code review happens: comments, requested changes, and automated checks
all attach to the PR, not to individual commits.

\begin{enumerate}
  \item Create a branch with a descriptive name:
        \code{git switch -c fix/uart-parity}.
  \item Commit your changes, then \code{git push -u origin
        fix/uart-parity}.
  \item On \GitHub{}, open a pull request from that branch into the
        target branch (usually \code{main}).
  \item Address review comments by committing more changes to the same
        branch and pushing again; the PR updates automatically.
  \item Once approved, merge it (or let a maintainer merge it) and
        delete the branch.
\end{enumerate}

\begin{ruleofthumb}
One pull request, one concern. A PR that fixes a bug, reformats
unrelated code and adds a new feature is nearly impossible to review
properly.
\end{ruleofthumb}

\begin{gotcha}{Diverged history after a rewritten branch}
If a branch's history was rewritten upstream (a rebase, an
\code{--amend} that was already pushed), a plain \code{git pull} will
refuse or create a tangled merge. \code{git pull --rebase} replays your
local commits on top of the new history instead; ask before force-pushing
a shared branch, since it can overwrite other people's work.
\end{gotcha}

\section{Issues and project boards}

\textbf{Issues} track bugs, tasks and questions independently of any
particular commit. A commit message or PR description containing
\code{Fixes \#12} automatically closes issue 12 when that commit lands
on the default branch. Labels, milestones and project boards organise
issues but are \GitHub{} features layered on top of \Git{} itself, not
part of the repository history.

\section{A minimal README}

Every repository should have a top-level \file{README.md}, since
\GitHub{} renders it automatically on the repository's front page.

\begin{txtcode}
# Project name

One paragraph: what this is and why it exists.

## Build
    make all

## Layout
    src/     source files
    test/    testbenches / unit tests
\end{txtcode}

\section{Checklist}
\label{sec:github:checklist}

\begin{itemize}
  \item \code{origin} is the remote you cloned from; add \code{upstream}
        when working from a fork.
  \item Use SSH remotes for repositories you push to often.
  \item One branch, one pull request, one concern.
  \item Never force-push a branch other people are also working on
        without asking first.
  \item A short \file{README.md} and a \file{.gitignore} are the two
        files every repository should start with.
\end{itemize}
